
\section{A reflection-symmetry barrier beyond axisymmetry}
\label{sec:xxxvi-reflection}

We seek a genuinely nonaxisymmetric, nonzero-initial-data,
exactly unforced comparison which retains the positive azimuthal
velocity of the original \emph{forced} Section~10 source
\cite{openai2026} at its distinguished shrinking trajectory.
The earlier circulation and averaged-torque estimates do not
distinguish angular oscillations that vanish at that trajectory.
Here we identify an additional, elementary but decisive discrete
symmetry barrier. It applies irrespective of the initial critical
norm and is verified on an explicitly prescribed large-data,
globally smooth nonaxisymmetric comparison. The symmetry principle
is classical; its specialization to the source trajectory is
the result of this section.

Put $R=\operatorname{diag}(1,-1,1)$ and introduce the
orthogonal reflection operator
\[
(\mathscr Ru)(x,y,z)=R\,u(x,-y,z),\qquad
u^{+}=\tfrac12(u+\mathscr Ru),\quad
u^{-}=\tfrac12(u-\mathscr Ru).
\]
The superscripts denote \emph{reflection parity}, not
Fourier signs or the axial no-swirl condition.

\begin{theorem}[Exact reflection barrier for the published shrinking-core swirl]
\label{thm:xxxvi-reflection}
Let $u$ be any smooth finite-energy, exactly unforced
Navier--Stokes solution on $\mathbb R^3\times[0,T)$,
whose smooth divergence-free initial datum is invariant
under $\mathscr R$. Then $u=\mathscr Ru$ for as long as
the smooth solution exists. In particular, for every $r>0$,
$z$ and $t<T$,
\begin{equation}\label{eq:xxxvi-reflection-node}
\boxed{\quad u_\theta(r,0,z,t)=u_y(r,0,z,t)=0,\quad}
\end{equation}
and its relative-to-$z$ angular momentum
$\Gamma=ru_\theta=xu_y-yu_x$ satisfies
\[
\boxed{\quad
\bar\Gamma=0,\qquad
\overline{u'_r\Gamma'}=\overline{u'_z\Gamma'}=0,
\qquad \mathcal T=0,\quad}
\]
where $\mathcal T$ is the exact nonaxisymmetric
mean-transport source of Edition~XXXV.
Such a flow cannot retain the original forced
source's positive azimuthal amplitude along
$x_\tau=(\sqrt{2X_{\rm in}\tau},0,0)$.
This holds for arbitrarily large smooth initial data
with the stated symmetry.
\end{theorem}
\begin{proof}
The Navier--Stokes equations and Leray projection
commute with simultaneous spatial/vector reflection
$\mathscr R$, and pressure transforms as
$p(x,y,z)\mapsto p(x,-y,z)$.
Smooth finite-energy uniqueness implies that
$\mathscr Ru(t)$ and $u(t)$ agree when their initial
data agree. On $y=0$, vector invariance forces $u_y=0$.
Since $e_\theta=e_y$ at $x=r>0,y=0$, this is
\eqref{eq:xxxvi-reflection-node}.
In cylindrical coordinates, reflection maps
$\theta\mapsto-\theta$, with $u_r,u_z$ even
and $u_\theta,\Gamma$ odd.
Every angular mean of an odd function vanishes;
$u_r\Gamma$ and $u_z\Gamma$ are odd as well.
The exact flux formula of Edition~XXXV therefore
has both covariance terms zero and $\mathcal T=0$.
The published source has positive
$u_{B,\theta}(x_\tau,1-\tau)\sim
e_0\tau^{-1/2-h}$, incompatible with the zero trace.
\end{proof}

\begin{theorem}[Evolution equation and quantitative necessary parity breaking]
\label{thm:xxxvi-odd-transport}
For a general smooth unforced solution $u$ (without
any reflection assumption), its parity-odd part obeys
the \emph{exact} projected linearized equation
\begin{equation}\label{eq:xxxvi-odd-pde}
\boxed{\quad
\partial_tu^{-}-\nu\Delta u^{-}
+\mathbb P\big[(u^{+}\cdot\nabla)u^{-}
+(u^{-}\cdot\nabla)u^{+}\big]=0.
\quad}
\end{equation}
If $q=\|u^{-}(0)\|_2$ and
$J(t)=\int_0^t\|\nabla u^{+}(s)\|_\infty\,ds$,
then
\begin{equation}\label{eq:xxxvi-odd-energy}
\boxed{\quad\|u^{-}(t)\|_2\le q e^{J(t)}.\quad}
\end{equation}
Suppose at the original positive-$x$ source trajectory
\[
|u_\theta(x_\tau,1-\tau)|
\ge c\,\tau^{-A},\quad
c>0,\quad A=1/2+h,\quad h>0.
\]
Then necessarily $q>0$, and at every sufficiently
small $\tau>0$ the parity-odd velocity satisfies
\begin{align}
\boxed{\quad
\|\nabla u^{-}(1-\tau)\|_\infty
&\ge\left(\frac{\pi c^5}
{24q^2e^{2J(1-\tau)}}\right)^{1/3}
\tau^{-5A/3},\quad}
\label{eq:xxxvi-odd-gradient}\\
\boxed{\quad
\|\Delta u^{-}(1-\tau)\|_2
&\ge\left(\frac{c}
{C_2q^{1/4}e^{J(1-\tau)/4}}\right)^{4/3}
\tau^{-4A/3},\quad}
\label{eq:xxxvi-odd-h2}
\end{align}
where $C_2=4/(\pi\sqrt2\,3^{7/8})$ is an
admissible explicit whole-space interpolation constant.
For the published swirl-retention test one can set
$c=e_0/4$. These are \emph{necessary}
quantitative costs, not existence estimates.
\end{theorem}
\begin{proof}
Write $u=u^{+}+u^{-}$,
$\mathscr Ru=u^{+}-u^{-}$.
Subtract the two projected unforced equations
and divide by two. The quadratic terms
$(u^+\cdot\nabla)u^+$ and
$(u^-\cdot\nabla)u^-$ cancel, yielding
\eqref{eq:xxxvi-odd-pde}.
Both parts remain divergence-free.
Its $L^2$ pairing with $u^-$ eliminates
$(u^+\cdot\nabla)u^-$ and controls the other
advection by $\|\nabla u^{+}\|_\infty$.
Gronwall proves \eqref{eq:xxxvi-odd-energy}.
On the plane $y=0$, the parity-even part has
zero $y$-component and $(u^{-})_y=u_y=u_\theta$.
Hence $\|u^{-}(1-\tau)\|_\infty\ge c\tau^{-A}$.
If $q=0$, the Gronwall bound implies $u^-=0$,
contradicting this inequality. For $q>0$ apply
the two whole-space interpolation estimates
of Edition~XXXIII to $u^-$,
using its $L^2$ bound $qe^J$:
\[
\|u^-\|_\infty^5\le
(24/\pi)\|u^-\|_2^2\|\nabla u^-\|_\infty^3,\qquad
\|u^-\|_\infty\le
C_2\|u^-\|_2^{1/4}\|\Delta u^-\|_2^{3/4}.
\]
Rearranging gives the asserted lower bounds.
\end{proof}

\begin{theorem}[Explicit large-data nonaxisymmetric global regular control]
\label{thm:xxxvi-shifted}
Fix $\nu>0$ and $d=1$, and set
\[
\boxed{
v_0(x,y,z)=2\nu e^{-((x-d)^2+y^2+z^2)}
\bigl((x-d)z,\ yz,\ 1-(x-d)^2-y^2\bigr).
}
\]
This is a smooth, Schwartz, divergence-free
initial velocity that is \emph{nonaxisymmetric}
relative to the source's $z$ axis.
It has the same exact large critical Fourier norm
as Edition~XXXV, since translation multiplies the
Fourier transform by a unit-modulus phase:
\[
\frac{(2\pi)^{-3/2}}{\nu}
\|v_0\|_{\mathcal X^{-1}}
=\frac{\sqrt\pi}{2}>\frac1{16}.
\]
The global axisymmetric-no-swirl existence theorem
applies around the displaced parallel axis $x=d,y=0$.
Its unique smooth exactly unforced solution
satisfies the reflection invariance of
Theorem~\ref{thm:xxxvi-reflection} for all time.
At $t=0$, its relative-to-source-axis angular momentum is
\begin{equation}\label{eq:xxxvi-shifted-gamma}
\boxed{\quad
\Gamma_0=2\nu d\,r z\,
e^{-(r^2+d^2-2dr\cos\theta+z^2)}\sin\theta.\quad}
\end{equation}
In particular, at $(r,z,\theta)=(1,1,\pi/2)$,
$\Gamma_0=2\nu e^{-3}\ne0$, and the sharp angular
trace estimate implies
\[
\|\partial_\theta\Gamma_0(1,\cdot,1)\|_{L^2_\theta}
\ge 2\nu e^{-3}\sqrt{6/\pi}>0.
\]
Nevertheless $\mathcal T=0$ identically throughout
the exact unforced evolution, and
$v_\theta(r,0,z,t)=0$. This explicitly refutes
the inference that large data, nonaxisymmetry and
nonzero angular oscillation suffice for growth
$\sim\tau^{-h}$ or retention of the source singularity.
\end{theorem}
\begin{proof}
Translate the globally regular no-swirl solution
of Edition~XXXV by the fixed vector $(d,0,0)$.
Navier--Stokes is translation invariant; its pressure
translates as well, with zero external forcing.
Its global smoothness and exact critical norm are unchanged.
The displayed datum is reflection invariant across
$y=0$, so Theorem~\ref{thm:xxxvi-reflection}
applies for all time. Direct Cartesian algebra gives
$\operatorname{div}v_0=0$ and
$\Gamma_0=x(v_0)_y-y(v_0)_x
=2\nu d y z e^{-((x-d)^2+y^2+z^2)}$,
which is \eqref{eq:xxxvi-shifted-gamma}.
The angular trace bound and torque conclusion
follow from its odd reflection parity.
\end{proof}

\paragraph{Implication for the full zero-force construction.}
A candidate retaining the original source's
positive swirl \emph{must} break $y$-reflection
already in its initial data. Its parity-odd
component must subsequently meet the
quantitative bounds \eqref{eq:xxxvi-odd-gradient}
and \eqref{eq:xxxvi-odd-h2} (or its parity-even
strain integral must grow sufficiently fast).
Such symmetry breaking is necessary but
not sufficient for unforced blowup.
The translated Gaussian control is an
actual zero-force whole-space solution,
not merely a kinematic stress or velocity.
No source-matched singular nonaxisymmetric
correction or end-to-end Lean PDE
formalization has been obtained.
